%%%%%%%%%%%%%%%%%%%%%%%%%%%%%%%%%%%%%%%%%%%%%%%%%%%%%%%%%%%%%%%%%%%%%%%%%%%%%
%% Contact Maps via S-Entropy Bidirectional Dijkstra:
%% A Theory of Individuation by Negation in Bounded Resolvable Spaces
%%
%% Author: Kundai Farai Sachikonye
%%         Technical University of Munich
%%%%%%%%%%%%%%%%%%%%%%%%%%%%%%%%%%%%%%%%%%%%%%%%%%%%%%%%%%%%%%%%%%%%%%%%%%%%%

\documentclass[twocolumn]{article}

\usepackage[a4paper, top=2.5cm, bottom=2.5cm, left=2cm, right=2cm,
            columnsep=0.8cm]{geometry}
\usepackage{amsmath}
\usepackage{amsthm}
\usepackage{amssymb}
\usepackage{amsfonts}
\usepackage{mathtools}
\usepackage[ruled,vlined,linesnumbered]{algorithm2e}
\usepackage[colorlinks=true, linkcolor=blue, citecolor=blue, urlcolor=blue]{hyperref}
\usepackage[backend=bibtex, style=numeric, sorting=none]{biblatex}
\addbibresource{references.bib}
\usepackage{microtype}
\usepackage{graphicx}
\usepackage{booktabs}
\usepackage{enumitem}
\usepackage{xcolor}
\usepackage{thmtools}

%% Theorem environments
\theoremstyle{plain}
\newtheorem{theorem}{Theorem}[section]
\newtheorem{lemma}[theorem]{Lemma}
\newtheorem{proposition}[theorem]{Proposition}
\newtheorem{corollary}[theorem]{Corollary}

\theoremstyle{definition}
\newtheorem{definition}[theorem]{Definition}
\newtheorem{example}[theorem]{Example}
\newtheorem{remark}[theorem]{Remark}
\newtheorem{assumption}[theorem]{Assumption}

%% Notation macros
\newcommand{\Omega}{\Omega}
\newcommand{\Borel}{\mathcal{B}}
\newcommand{\PP}{\mathcal{P}}
\newcommand{\CM}{\mathcal{CM}}
\newcommand{\NN}{\mathcal{N}}
\newcommand{\TC}{\mathcal{T}_C}
\newcommand{\GC}{G_C}
\newcommand{\SS}{\mathcal{S}}
\newcommand{\HH}{\mathcal{H}}
\newcommand{\Sk}{S_k}
\newcommand{\St}{S_t}
\newcommand{\Se}{S_e}
\newcommand{\bS}{\mathbf{S}}
\newcommand{\dS}{d_{\mathcal{S}}}
\newcommand{\partdep}{\beta_{\PP}}
\newcommand{\bstar}{\beta_*}
\newcommand{\mumin}{\mu_{\min}}
\newcommand{\reals}{\mathbb{R}}
\newcommand{\naturals}{\mathbb{N}}
\newcommand{\eps}{\varepsilon}
\newcommand{\qed}{\hfill$\blacksquare$}
\newcommand{\norm}[1]{\left\lVert #1 \right\rVert}
\newcommand{\abs}[1]{\left\lvert #1 \right\rvert}
\newcommand{\interior}[1]{\mathring{#1}}
\newcommand{\closure}[1]{\overline{#1}}

\title{Contact Maps via S-Entropy Bidirectional Dijkstra:\\
A Theory of Individuation by Negation in Bounded Resolvable Spaces}

\author{Kundai Farai Sachikonye\\
\textit{Technical University of Munich}\\
\texttt{k.sachikonye@tum.de}}

\date{\today}

\begin{document}

\maketitle

%%%%%%%%%%%%%%%%%%%%%%%%%%%%%%%%%%%%%%%%%%%%%%%%%%%%%%%%%%%%%%%%%%%%%%%%%%%%%
\begin{abstract}
%%%%%%%%%%%%%%%%%%%%%%%%%%%%%%%%%%%%%%%%%%%%%%%%%%%%%%%%%%%%%%%%%%%%%%%%%%%%%
We develop a rigorous mathematical framework for image segmentation and
holographic reconstruction grounded in the principle of \emph{individuation by
negation}: a region is individuated not by its intrinsic properties alone, but
by the totality of what it is not. This deceptively simple principle has
profound algorithmic consequences. In a bounded resolvable space
$(\Omega, d, \mu)$ --- a compact metric space equipped with a Borel measure and
a positive resolution floor $\mumin > 0$ --- we prove that individuation is
intrinsically non-instantaneous: every act of distinguishing a region $A$ from
its complement requires at least one committed partition step of cost $\geq
\mumin$. This non-instantaneity forbids zero-cost boundary detection and
provides the foundational justification for finite-resolution algorithms.

Central to our theory is the concept of \emph{contact} between regions: two
regions $A$ and $B$ are in contact if and only if their topological boundaries
share at least one point at the minimum achievable partition depth $\bstar$.
We prove that contact is an invariant of the partition topology --- once
established, it persists under refinement --- and that it is \emph{irreducible}:
no strictly smaller topological fact carries equivalent information across all
depth levels.

To compute contact maps efficiently, we introduce \emph{S-Entropy Bidirectional
Dijkstra} (SEBD), an algorithm that operates on a three-dimensional state space
of normalized entropy coordinates $\bS(A) = (\Sk(A), \St(A), \Se(A)) \in
[0,1]^3$ capturing kinetic, thermal, and energetic entropy of each image region.
SEBD employs \emph{virtual substates} --- mathematically necessary intermediary
states that may lie outside the physical state space $[0,1]^3$ --- to unfold
the backward search frontier, enabling exact computation of the minimum-cost
path in $O(|\mathcal{R}|^{3/2})$ time where $|\mathcal{R}|$ is the number of
distinguishable regions.

The contact map drives a \emph{contact-driven holographic slicing} algorithm
that produces a 3D holographic representation of the image. Crucially, the
holographic z-axis is not a physical dimension nor a predetermined parameter: it
is the ordering of contact resolutions by SEBD cost, making depth an emergent
property of the image's partition topology. We prove completeness of the slicing
algorithm, faithful representation of partition structure, and residue
propagation of new contacts at each resolution step. Numerical verification on
synthetic and natural image examples is provided, together with applications to
medical image segmentation, 3D holographic reconstruction, image compression,
and scene understanding.
\end{abstract}

%%%%%%%%%%%%%%%%%%%%%%%%%%%%%%%%%%%%%%%%%%%%%%%%%%%%%%%%%%%%%%%%%%%%%%%%%%%%%
\section{Introduction}
\label{sec:introduction}
%%%%%%%%%%%%%%%%%%%%%%%%%%%%%%%%%%%%%%%%%%%%%%%%%%%%%%%%%%%%%%%%%%%%%%%%%%%%%

The problem of partitioning an image into meaningful regions is perhaps the
oldest open problem in computer vision, and yet the dominant paradigm --- detect
edges by finding intensity discontinuities, then segment --- suffers from a
fundamental conceptual confusion: it conflates \emph{proximity} with
\emph{contact}, and \emph{intensity discontinuity} with \emph{partition
boundary}.

Consider two adjacent gray regions in an image, both of uniform intensity but
separated by a thin dark band. A Canny edge detector \cite{canny1986computational}
reports two edges (one on each side of the dark band) and treats the dark band
as a third region. The contact structure among the three regions is entirely
obscured. Now consider the same image but with the dark band infinitesimally
thin: in the limit, the detector reports two coincident edges, yet the two gray
regions are still distinct objects. The edge detector cannot distinguish
``two adjacent objects'' from ``one object with a dark stripe.'' This failure is
not a limitation of the specific algorithm; it is a consequence of the
philosophical confusion between detecting intensity gradients and resolving
partition boundaries.

The classical watershed transform \cite{beucher1979use, vincent1991watersheds}
partially addresses this by treating the image as a topographic surface and
flooding it from local minima. The watershed boundaries are contours, but the
flooding order is determined by intensity values rather than by the logical
structure of the partition. The normalized cuts approach
\cite{shi2000normalized} avoids intensity-gradient traps by optimizing a global
graph criterion, but the criterion is not derived from any topological theory of
what a boundary \emph{is}. Graph-based methods
\cite{felzenszwalb2004efficient, boykov2004experimental} offer computational
efficiency at the cost of theoretical clarity about what makes two regions
``separate.''

\subsection{Individuation by Negation}

The present paper proposes an entirely different starting point. A region is not
defined by what it contains; it is defined by what it \emph{excludes}. Formally:

\begin{center}
\emph{$A$ is individuated from $\Omega$ by negating everything $A$ is not.}
\end{center}

This principle --- which we call \emph{individuation by negation} --- has roots
in classical philosophy (the Spinozan \emph{omnis determinatio est negatio}),
but here we give it a precise mathematical formulation in the language of measure
theory and topology. The key insight is that this negation is \emph{not free}: it
costs at least $\mumin > 0$ in measure, because in a bounded resolvable space,
every separator between two regions must have positive measure. There is no
such thing as an infinitely thin boundary in a finite-resolution world.

This insight immediately implies that individuation is a process, not an event:
it takes at least one committed, irreversible partition step to separate $A$ from
its complement. Each such step creates a \emph{residue} --- new boundary between
$A$ and its other neighbors --- which must itself be resolved, propagating the
individuation process through the entire contact graph.

\subsection{Contact as the Fundamental Invariant}

Given the non-instantaneity of individuation, the natural question is: what
information must be preserved and propagated from one partition step to the next?
We prove that the answer is \emph{contact}: whether two regions share a boundary
point. Contact is invariant under partition refinement (once two regions touch,
they always touch at higher resolution) and irreducible (no coarser topological
fact carries the same information). The contact relation generates a topology
$\TC$ on the set of distinguishable regions, and the contact graph $\GC$ is the
combinatorial skeleton of this topology.

\subsection{S-Entropy State Space and SEBD}

To quantify the ``distance'' between two regions in the contact map, we introduce
the \emph{S-entropy state space} $\SS = [0,1]^3$, whose three axes are
normalized measures of kinetic entropy ($\Sk$: gradient magnitude density),
thermal entropy ($\St$: intensity distribution entropy), and energetic entropy
($\Se$: contrast ratio). The S-entropy coordinates of a region are computable
directly from the image data and are provably well-defined for all distinguishable
regions.

The contact map assigns to each adjacent pair $(A,B)$ the Euclidean distance
$\dS(\bS(A), \bS(B))$ in this state space. To compute shortest paths in the
contact map efficiently, we introduce \emph{S-Entropy Bidirectional Dijkstra}
(SEBD), which runs forward search from $\bS(A)$ and backward search from $\bS(B)$
using \emph{virtual substates}: off-shell states in $\reals^3 \setminus [0,1]^3$
that represent mathematically necessary intermediaries in the path decomposition.

\subsection{Holographic Slicing}

Finally, the contact map drives a \emph{contact-driven holographic slicing}
algorithm. At each step, the pair $(A,B)$ with minimum contact cost is resolved
--- $A$ and $B$ are individuated --- and the boundary $\partial A \cap \partial B$
is recorded as a ``slice'' at holographic depth $z = \CM(A,B)$. This produces
a 3D holographic representation of the image whose z-axis is the ordering of
contact resolutions. The depth is not physical; it is the topology of the image
made spatial.

\subsection{Contributions}

Our main contributions are:
\begin{enumerate}[label=(\roman*)]
\item A rigorous mathematical framework for bounded resolvable spaces and the
  principle of individuation by negation (Sections~\ref{sec:setting}--\ref{sec:individuation}).
\item The theory of contact as a topological invariant, including proofs of
  Contact Invariance and Contact Irreducibility (Section~\ref{sec:contact}).
\item The S-entropy state space formalism for image regions
  (Section~\ref{sec:resolution}).
\item The S-Entropy Contact Map and the SEBD algorithm with proofs of correctness
  and complexity (Section~\ref{sec:contactmap}).
\item The contact-driven holographic slicing algorithm with proofs of
  Completeness, Residue Propagation, and Hologram Faithfulness
  (Sections~\ref{sec:slicing}--\ref{sec:properties}).
\item Numerical verification on worked examples and qualitative experimental
  analysis (Section~\ref{sec:numerical}).
\item Applications to segmentation, 3D reconstruction, compression, and scene
  understanding (Section~\ref{sec:applications}).
\end{enumerate}

\subsection{Organization}

The paper is organized as follows. Section~\ref{sec:setting} establishes the
mathematical setting. Section~\ref{sec:individuation} develops individuation by
negation. Section~\ref{sec:contact} establishes contact as a topological
invariant. Section~\ref{sec:resolution} specializes to image spaces and defines
the S-entropy coordinates. Section~\ref{sec:contactmap} defines the contact map
and SEBD algorithm. Section~\ref{sec:slicing} presents contact-driven holographic
slicing. Section~\ref{sec:properties} collects and proves all main theoretical
results. Section~\ref{sec:numerical} provides numerical verification.
Section~\ref{sec:applications} discusses applications. Section~\ref{sec:conclusion}
concludes.

%%%%%%%%%%%%%%%%%%%%%%%%%%%%%%%%%%%%%%%%%%%%%%%%%%%%%%%%%%%%%%%%%%%%%%%%%%%%%
\section{Mathematical Setting}
\label{sec:setting}
%%%%%%%%%%%%%%%%%%%%%%%%%%%%%%%%%%%%%%%%%%%%%%%%%%%%%%%%%%%%%%%%%%%%%%%%%%%%%

We establish the foundational mathematical framework. Throughout, we work in a
setting that is general enough to encompass digital images, discrete graphs, and
continuous measure spaces, while specific enough to make the resolution floor
a theorem rather than an assumption.

\begin{definition}[Bounded Resolvable Space]
\label{def:brs}
A \emph{bounded resolvable space} is a triple $(\Omega, d, \mu)$ where:
\begin{enumerate}[label=(\alph*)]
\item $\Omega$ is a compact metric space with metric $d$;
\item $\mu$ is a non-negative, non-atomic Borel measure on $\Omega$ with
  $\mu(\Omega) < \infty$;
\item there exists a \emph{resolution constant} $\delta > 0$ such that for every
  open ball $B(x, r)$ with $r < \delta$, the measure $\mu(B(x,r))$ is bounded
  away from zero: $\mu(B(x,r)) \geq c_0 r^n$ for some constants $c_0 > 0$ and
  $n \geq 1$ (the effective dimension of the space).
\end{enumerate}
\end{definition}

\begin{remark}
The non-atomicity of $\mu$ ensures that singletons have zero measure, so that
the boundary of a region always has strictly smaller measure than the interior.
The lower density bound in condition (c) is a mild regularity assumption satisfied
by Lebesgue measure on $\reals^n$, by the normalized counting measure on finite
grids, and by the Hausdorff measure on regular fractals.
\end{remark}

\begin{definition}[Partition]
\label{def:partition}
A \emph{partition} of $\Omega$ is a finite collection $\PP = \{A_1, \ldots, A_k\}$
of measurable subsets of $\Omega$ such that:
\begin{enumerate}[label=(\alph*)]
\item $A_i \cap A_j = \emptyset$ for all $i \neq j$ (mutual disjointness);
\item $\mu\!\left(\Omega \setminus \bigcup_{i=1}^k A_i\right) = 0$ (essential coverage).
\end{enumerate}
We say that $\PP$ is a \emph{valid partition} if additionally each $A_i$ is
connected (in the topological sense) and $\mu(A_i) > 0$ for all $i$.
\end{definition}

\begin{remark}
The essential coverage condition allows for a zero-measure set of ``boundary
pixels'' that are not assigned to any region. This is natural in the digital
setting, where pixel centers belong to regions but pixel edges belong to no
region.
\end{remark}

\begin{definition}[Separator]
\label{def:separator}
Let $\PP = \{A_1, \ldots, A_k\}$ be a partition of $\Omega$ and let $A = A_i$
for some $i$. A \emph{separator} of $A$ in $\PP$ is any $B = A_j \in \PP$
with $j \neq i$ such that $B$ is adjacent to $A$, i.e., such that
$\partial A \cap \partial B \neq \emptyset$ in the metric topology of $\Omega$.
We denote the set of all separators of $A$ in $\PP$ by $\mathrm{Sep}(A, \PP)$.
\end{definition}

\begin{definition}[Partition Depth]
\label{def:partdep}
Let $\PP$ be a partition and $A \in \PP$. The \emph{partition depth} of $A$ in
$\PP$ is
\[
  \partdep(A) \;=\; \inf_{B \in \mathrm{Sep}(A,\PP)} \mu(B).
\]
When $\mathrm{Sep}(A,\PP) = \emptyset$ (i.e., $A = \Omega$), we set
$\partdep(A) = +\infty$ by convention.
\end{definition}

\begin{definition}[Resolution Floor]
\label{def:resfloor}
A bounded resolvable space $(\Omega, d, \mu)$ has \emph{resolution floor}
$\mumin > 0$ if for every valid partition $\PP$ and every $A \in \PP$,
every separator $B \in \mathrm{Sep}(A,\PP)$ satisfies $\mu(B) \geq \mumin$.
Equivalently, $\partdep(A) \geq \mumin$ for all $A$ and all valid $\PP$.
\end{definition}

\begin{definition}[Minimum Partition Depth]
\label{def:minpartdep}
The \emph{minimum partition depth} of the space is
\[
  \bstar \;=\; \inf_{\PP \text{ valid}} \inf_{A \in \PP} \partdep(A).
\]
\end{definition}

The key structural result of this section is:

\begin{lemma}[Positivity of Minimum Partition Depth]
\label{lem:posfloor}
In any bounded resolvable space with resolution floor $\mumin > 0$, the minimum
partition depth satisfies $\bstar \geq \mumin > 0$.
\end{lemma}

\begin{proof}
We must show that for every valid partition $\PP$ and every $A \in \PP$, the
partition depth $\partdep(A) \geq \mumin$.

By Definition~\ref{def:resfloor}, for every valid partition $\PP$ and every $A \in \PP$,
every separator $B \in \mathrm{Sep}(A, \PP)$ satisfies $\mu(B) \geq \mumin$.
Therefore
\[
  \partdep(A) \;=\; \inf_{B \in \mathrm{Sep}(A,\PP)} \mu(B) \;\geq\; \mumin \;>\; 0.
\]
Since this holds for all valid $\PP$ and all $A \in \PP$, we conclude
\[
  \bstar \;=\; \inf_{\PP} \inf_{A \in \PP} \partdep(A) \;\geq\; \mumin \;>\; 0.
\]

To see that the resolution floor assumption is not vacuous, we verify it for
the digital image setting. Let $\Omega = \{1,\ldots,M\} \times \{1,\ldots,N\}$
be a finite grid, $d$ the Chebyshev metric, and $\mu$ the normalized counting
measure $\mu(S) = |S|/(MN)$. Then any non-empty set $B \subset \Omega$ satisfies
$\mu(B) \geq 1/(MN)$. Since every separator is non-empty, we have $\mumin \geq
1/(MN) > 0$. In the continuous case $\Omega \subset \reals^2$ with Lebesgue
measure, we impose a minimum region size requirement: a ``region'' must contain
a ball of radius at least $\delta/2$, so $\mu(A) \geq c_0 (\delta/2)^2 > 0$ for
all separators, giving $\mumin = c_0 \delta^2 / 4 > 0$.

Thus in all cases of practical interest, $\bstar \geq \mumin > 0$. \qed
\end{proof}

\begin{remark}
Lemma~\ref{lem:posfloor} is the single most important fact in this paper. It says:
\emph{there are no infinitely thin boundaries in a bounded resolvable space}.
Every separator between two regions has positive measure. This is why individuation
cannot be instantaneous (Theorem~\ref{thm:noninstant}), why contact is well-defined
(Definition~\ref{def:contact}), and why the SEBD algorithm terminates in finite time
(Theorem~\ref{thm:complexity}).
\end{remark}

%%%%%%%%%%%%%%%%%%%%%%%%%%%%%%%%%%%%%%%%%%%%%%%%%%%%%%%%%%%%%%%%%%%%%%%%%%%%%
\section{Individuation by Negation}
\label{sec:individuation}
%%%%%%%%%%%%%%%%%%%%%%%%%%%%%%%%%%%%%%%%%%%%%%%%%%%%%%%%%%%%%%%%%%%%%%%%%%%%%

The central philosophical and mathematical concept of this paper is
\emph{individuation by negation}: the idea that a region acquires its identity
not through its own positive properties, but through the systematic negation
of everything else in the space.

\begin{definition}[Complement and Negation]
\label{def:negation}
Let $\PP = \{A_1, \ldots, A_k\}$ be a partition of $\Omega$ and let $A = A_i$.
The \emph{negation} of $A$ in $\PP$ is
\[
  \NN(A) \;=\; \bigcup_{j \neq i} A_j.
\]
The \emph{individuation} of $A$ by negation is the identity
\[
  A \;=\; \Omega \setminus \NN(A)
\]
(equality up to a set of $\mu$-measure zero).
\end{definition}

\begin{remark}
The content of Definition~\ref{def:negation} is not the set-theoretic tautology
$A = \Omega \setminus (\Omega \setminus A)$. The content is that $A$ is
\emph{determined} by the collection of all other regions, not just by its own
intrinsic description. If we remove any one of the $A_j$ ($j \neq i$) from
$\NN(A)$, the identity $A = \Omega \setminus \NN(A)$ fails: the right-hand side
now contains the removed region. Thus $A$ depends on the \emph{totality} of its
negation, not merely on its immediate neighbors.
\end{remark}

\begin{definition}[Committed Partition Step]
\label{def:committed}
A \emph{committed partition step} is an operation that takes a partition $\PP$
and produces a refined partition $\PP' \succ \PP$ (meaning $\PP'$ is obtained by
splitting some element of $\PP$) such that the boundary measure of the new
separator is at least $\mumin$:
\[
  \mu(\partial A' \cap \partial B') \;\geq\; \mumin
\]
for the newly created pair $(A', B')$. We say the step is \emph{irreversible}
because once the separator has been committed to at cost $\geq \mumin$, it cannot
be undone without violating the resolution floor.
\end{definition}

\begin{definition}[Individuation Residue]
\label{def:residue}
Let $A$ and $B$ be two adjacent regions in a partition $\PP$, and suppose we
individuate $A$ from $B$ by a committed partition step that separates them.
The \emph{residue} of this individuation is
\[
  R(A,B) \;=\; \partial A \cap \partial B,
\]
the shared boundary that is revealed by the separation. This residue is not
destroyed by the individuation step; rather, it becomes the locus of new
potential contact relationships between $A$ and $B$'s other neighbors. More
precisely, for any $C \in \PP \setminus \{A, B\}$ that was previously adjacent
only to $B$ (not to $A$), the residue may create a new adjacency between $A$
and $C$ in the refined partition $\PP'$.
\end{definition}

The main theorem of this section formalizes the non-instantaneity of individuation.

\begin{theorem}[Non-Instantaneity of Individuation]
\label{thm:noninstant}
Let $(\Omega, d, \mu)$ be a bounded resolvable space with resolution floor
$\mumin > 0$. Let $\PP$ be any valid partition and let $A \in \PP$ be any region
that is not equal to $\Omega$ (i.e., $A$ has at least one separator). Then the
individuation of $A$ in $\PP$ requires at least one committed partition step of
cost $\geq \mumin$.
\end{theorem}

\begin{proof}
We must show that any process of individuating $A$ --- producing the identity
$A = \Omega \setminus \NN(A)$ --- requires at least one committed partition step.

Suppose, for contradiction, that there exists an individuation of $A$ that
requires \emph{zero} committed partition steps. Then the negation $\NN(A) =
\bigcup_{j \neq i} A_j$ is assembled without any partition step of cost $\geq
\mumin$, meaning that for every separator $B \in \mathrm{Sep}(A,\PP)$, we have
$\mu(B) < \mumin$. But by Definition~\ref{def:resfloor}, the resolution floor
requires $\mu(B) \geq \mumin$ for all valid separators $B$. This is a
contradiction.

More constructively: since $A \neq \Omega$, there exists at least one separator
$B \in \mathrm{Sep}(A,\PP)$. By Lemma~\ref{lem:posfloor}, $\mu(B) \geq \mumin > 0$.
The separation of $A$ from $B$ is a partition step that assigns boundary measure
$\mu(\partial A \cap \partial B) > 0$. Since $\mu(\partial A \cap \partial B) \leq
\mu(B)$, and $\mu(B) \geq \mumin$, this step costs at least $\mumin$ in measure.

To see that this step is irreversible: once the partition $\PP$ has been committed
to (i.e., $A$ and $B$ have been assigned to separate cells), any reassignment
that merges $A$ and $B$ into a single cell would violate the partition's
disjointness requirement without going through the full cost of a new partition
step. Thus the step is committed and irreversible.

Finally, we observe that one such step is \emph{sufficient} (in principle): if
$\PP$ is already a valid partition and $A$ already appears as a cell, then the
individuation is complete in one step (the step of defining $\PP$ itself). The
theorem asserts a lower bound of at least one step, not an exact count.

We conclude that the individuation of any non-trivial region in a bounded
resolvable space requires at least one committed partition step of cost $\geq
\mumin > 0$. \qed
\end{proof}

\begin{remark}
The theorem has a direct algorithmic interpretation: any algorithm that claims
to individuate a region in zero time (e.g., by detecting an edge with a
derivative filter and declaring the region boundaries to coincide with the edge)
is making an implicit assumption that violates the resolution floor. Edge
detection is not individuation; it is a proxy for individuation that fails
precisely when the proxy is needed most (thin separators, low-contrast
boundaries, textured regions).
\end{remark}

\begin{proposition}[Residue Chain]
\label{prop:residuechain}
Let $\PP = \{A_1, \ldots, A_k\}$ be a valid partition. Let the individuation
sequence $A_{i_1}, A_{i_2}, \ldots, A_{i_m}$ be the order in which regions are
individuated (i.e., separated from all their neighbors). Then the residue of each
individuation step generates new contacts that must be resolved in subsequent
steps. The total number of committed partition steps is at least $k-1$ (one per
separation in the spanning tree of the contact graph $\GC$).
\end{proposition}

\begin{proof}
The contact graph $\GC$ (defined precisely in the next section) has $k$ vertices
(one per region) and is connected (because $\Omega$ is connected and each region
is connected). Any connected graph on $k$ vertices requires at least $k-1$ edges
to be connected; a spanning tree has exactly $k-1$ edges. Each edge of the
spanning tree corresponds to at least one committed partition step separating
the two endpoint regions. Since the steps are irreversible and each one resolves
exactly one contact in the spanning tree, the total count is at least $k-1$. \qed
\end{proof}

%%%%%%%%%%%%%%%%%%%%%%%%%%%%%%%%%%%%%%%%%%%%%%%%%%%%%%%%%%%%%%%%%%%%%%%%%%%%%
\section{Contact as Topological Invariant}
\label{sec:contact}
%%%%%%%%%%%%%%%%%%%%%%%%%%%%%%%%%%%%%%%%%%%%%%%%%%%%%%%%%%%%%%%%%%%%%%%%%%%%%

We now formalize the concept of contact between regions and prove its topological
invariance and irreducibility.

\begin{definition}[Contact]
\label{def:contact}
Let $(\Omega, d, \mu)$ be a bounded resolvable space with minimum partition depth
$\bstar > 0$. Two regions $A, B \in \PP$ (in some valid partition $\PP$) are in
\emph{contact} at threshold $\tau \geq \bstar$ if
\[
  \partial A \cap \partial B \;\neq\; \emptyset
\]
where $\partial$ denotes the topological boundary in $(\Omega, d)$. We write
$C_\tau(A,B) = 1$ if they are in contact at threshold $\tau$, and $C_\tau(A,B) = 0$
otherwise. At the minimum threshold, we write $C(A,B) = C_{\bstar}(A,B)$.
\end{definition}

\begin{definition}[Contact versus Closeness]
\label{def:closeness}
The \emph{distance} between two regions $A$ and $B$ is
\[
  d(A,B) \;=\; \inf_{x \in A, y \in B} d(x,y).
\]
We say $A$ and $B$ are \emph{close} at scale $\eps > 0$ if $d(A,B) < \eps$.
Closeness and contact are \emph{distinct} concepts: $d(A,B) < \eps$ does not
imply $C(A,B) = 1$, and $C(A,B) = 1$ does not imply $d(A,B)$ is small.
\end{definition}

\begin{remark}[Contact vs.\ Proximity]
Consider the partition $\PP = \{A, G, B\}$ of a one-dimensional interval
$\Omega = [0,1]$ with $A = [0, 0.4]$, $G = [0.4, 0.6]$ (a gap region), and
$B = [0.6, 1]$. Then $d(A,B) = 0.2$, which can be made arbitrarily small by
shrinking $G$, yet $C(A,B) = 0$ because $\partial A \cap \partial B = \emptyset$
(they share no boundary point in the partition $\PP$). Conversely, in the
partition $\PP' = \{A', B'\}$ with $A' = [0,0.5]$ and $B' = [0.5,1]$, we have
$d(A', B') = 0$ and $C(A',B') = 1$ (they share the boundary point $0.5$).
This shows that contact is a property of the \emph{partition}, not of the metric.
\end{remark}

\begin{definition}[Contact Graph]
\label{def:contactgraph}
Let $\PP = \{A_1, \ldots, A_k\}$ be a valid partition of $\Omega$. The
\emph{contact graph} $\GC = (V, E)$ is the undirected graph with:
\begin{itemize}
\item vertex set $V = \{A_1, \ldots, A_k\}$ (one vertex per region);
\item edge set $E = \{(A_i, A_j) : C(A_i, A_j) = 1, i \neq j\}$.
\end{itemize}
\end{definition}

\begin{definition}[Contact Topology]
\label{def:contacttop}
The \emph{contact topology} $\TC$ on $V$ is the topology generated by the
contact relation: the open sets of $\TC$ are the unions of cliques in $\GC$
(sets of pairwise-contacting regions). Equivalently, $U \subset V$ is open in
$\TC$ if for every $A \in U$ and every $B$ in contact with $A$, the pair
$(A,B)$ is ``resolvable'' in $U$ (meaning the contact between them is registered
in $U$).
\end{definition}

\begin{definition}[Bilateral Negation]
\label{def:bilateral}
We say that region $A$ \emph{bilaterally negates} region $B$, written $A \perp B$,
if:
\begin{enumerate}[label=(\roman*)]
\item $A \not\subset B$ and $B \not\subset A$ (neither contains the other);
\item $C(A,B) = 1$ (they are in contact).
\end{enumerate}
Bilateral negation is more primitive than the global negation $\NN(A) =
\Omega \setminus A$: it is a pairwise relationship between two regions that
are simultaneously individuating each other.
\end{definition}

\begin{definition}[Monotone Contact Record]
\label{def:monotone}
The \emph{monotone contact record} at threshold $\tau$ is the set
\[
  M_\tau \;=\; \{(A,B) \in V \times V : C_\tau(A,B) = 1\}.
\]
We say the contact record is \emph{monotone} if $M_\tau \subseteq M_{\tau'}$
whenever $\tau \leq \tau'$ (contacts are preserved under partition refinement).
\end{definition}

\begin{theorem}[Contact Invariance]
\label{thm:invariance}
Let $(\Omega, d, \mu)$ be a bounded resolvable space. The contact record is
monotone: if $C_\tau(A,B) = 1$ for some threshold $\tau \geq \bstar$, then
$C_{\tau'}(A,B) = 1$ for all $\tau' \geq \tau$.
\end{theorem}

\begin{proof}
Suppose $C_\tau(A,B) = 1$, meaning $\partial A \cap \partial B \neq \emptyset$
at threshold $\tau$. Let $p \in \partial A \cap \partial B$ be a contact point.
We must show that $p \in \partial A' \cap \partial B'$ at threshold $\tau'$,
where $A', B'$ are the images of $A, B$ in the refined partition at threshold
$\tau'$.

Since $\tau' \geq \tau$, the refined partition $\PP'$ at threshold $\tau'$ is
obtained from $\PP$ (the partition at threshold $\tau$) by possibly splitting
some of its cells further. We consider two cases:

\textbf{Case 1}: Neither $A$ nor $B$ is split in the refinement from $\PP$ to
$\PP'$. Then $A' = A$ and $B' = B$, so $\partial A' \cap \partial B' =
\partial A \cap \partial B \ni p$, and $C_{\tau'}(A,B) = 1$ trivially.

\textbf{Case 2}: One or both of $A$, $B$ are split. Say $A$ is split into
$A_1' \cup A_2'$. By the connectedness of $A$ and the fact that the split
creates an internal separator, the contact point $p$ lies on $\partial A$, which
is preserved in $\partial A_1'$ or $\partial A_2'$ (the piece of $A$ that still
borders $B$). More precisely: $p \in \partial A \cap \partial B$, and since any
split of $A$ only introduces a new internal boundary, the external boundary
$\partial A \cap \partial B$ is a subset of $\partial A_1' \cup \partial A_2'$.
Hence $p \in \partial A_j' \cap \partial B'$ for some $j \in \{1,2\}$, meaning
$C_{\tau'}(A_j', B) = 1$.

In either case, the contact relationship is preserved (though it may now be
between a sub-region $A_j'$ and $B$ rather than between $A$ and $B$). Since
we track contact at the level of the finest partition, the contact record
$M_{\tau'} \supseteq M_\tau$. This shows monotonicity of the contact record.

Formally: for any $(A,B) \in M_\tau$, the pair $(A,B)$ or a refined variant
$(A_j', B_{j'}')$ remains in $M_{\tau'}$, so $M_\tau \subseteq M_{\tau'}$. \qed
\end{proof}

\begin{theorem}[Contact Irreducibility]
\label{thm:irreducible}
Contact is the minimum sufficient topological invariant that must be propagated
between partition thresholds: no strictly coarser topological fact about the
partition carries the same information across all depth levels.
\end{theorem}

\begin{proof}
We prove this by showing that any topological fact $F(A,B)$ that is implied by
contact ($C(A,B) = 1 \Rightarrow F(A,B)$) and is preserved under refinement is
either equivalent to contact or is trivially true.

Let $F$ be any topological predicate on pairs of regions that is:
\begin{enumerate}[label=(\roman*)]
\item \emph{invariant}: $C(A,B) = 1$ implies $F(A,B)$;
\item \emph{monotone}: $F(A,B)$ at threshold $\tau$ implies $F(A,B)$ at $\tau' \geq \tau$;
\item \emph{non-trivial}: there exist distinct, non-contacting regions $A, B$ with $F(A,B) = 0$.
\end{enumerate}

We claim $F(A,B) \Leftrightarrow C(A,B)$.

The direction $C(A,B) \Rightarrow F(A,B)$ is condition (i).

For the converse, suppose $F(A,B) = 1$ but $C(A,B) = 0$, i.e., $\partial A \cap
\partial B = \emptyset$. Since $A$ and $B$ are disjoint (they belong to a
partition) and their boundaries do not intersect, there exists an open set $U$
in $\Omega$ that contains $\partial A$ but not $\partial B$, and similarly an
open set $V$ containing $\partial B$ but not $\partial A$, with $U \cap V = \emptyset$.
This means $A$ and $B$ are separated in the topological sense \cite{munkres2000topology}.

By the definition of the contact topology $\TC$, separated regions are in
different connected components of $\TC$. Any topological predicate that depends
only on the $\TC$-structure of the pair $(A,B)$ must therefore evaluate to $0$
when $A$ and $B$ are in different components --- the same value as $C(A,B) = 0$.

If $F$ is finer than contact (i.e., $F(A,B) = 0$ whenever $C(A,B) = 0$), then
$F$ is equivalent to contact on the class of non-trivial partitions.

If $F$ is coarser than contact (i.e., $F(A,B) = 1$ for some pair with $C(A,B) = 0$),
then $F$ fails to distinguish contacting from non-contacting pairs, and by
condition (iii) it is trivially true for all pairs --- contradicting non-triviality.

Therefore, the only monotone, non-trivial topological predicate that implies and
is implied by contact is contact itself, establishing its irreducibility. \qed
\end{proof}

%%%%%%%%%%%%%%%%%%%%%%%%%%%%%%%%%%%%%%%%%%%%%%%%%%%%%%%%%%%%%%%%%%%%%%%%%%%%%
\section{The Resolution Floor in Image Spaces}
\label{sec:resolution}
%%%%%%%%%%%%%%%%%%%%%%%%%%%%%%%%%%%%%%%%%%%%%%%%%%%%%%%%%%%%%%%%%%%%%%%%%%%%%

We now specialize the abstract framework to the concrete setting of digital images,
where the S-entropy coordinates provide a computable and faithful representation
of each region's entropic character.

\subsection{Images as Bounded Resolvable Spaces}

\begin{definition}[Image Space]
\label{def:imagespace}
An \emph{image} is a measurable, bounded function $I: \Omega \to \reals$ where
$\Omega \subset \reals^2$ is compact. We equip $\Omega$ with the Euclidean
metric $d$ and the Lebesgue measure $\mu$ (normalized so $\mu(\Omega) = 1$).
We assume $I \in L^1(\Omega) \cap L^2(\Omega)$ and that $\nabla I$ exists
$\mu$-almost everywhere (e.g., $I \in W^{1,2}(\Omega)$, a Sobolev space
\cite{evans2010partial}).
\end{definition}

\begin{remark}
For digital images, $\Omega = \{1,\ldots,M\} \times \{1,\ldots,N\}$ is a finite
grid, $I$ is defined on grid points, and $\nabla I$ is replaced by the discrete
gradient (finite differences). All results below carry over to the discrete
setting by replacing integrals with sums.
\end{remark}

\begin{definition}[S-Entropy Coordinates]
\label{def:sentropy}
Let $A \subset \Omega$ be a measurable region with $\mu(A) > 0$. The
\emph{S-entropy coordinates} of $A$ are the triple $\bS(A) = (\Sk(A), \St(A), \Se(A))$
where:

\paragraph{Kinetic Entropy $\Sk(A)$} (normalized gradient magnitude density):
\[
  \Sk(A) \;=\; \frac{\displaystyle\frac{1}{\mu(A)} \int_A \norm{\nabla I}_2 \, d\mu}
                    {\displaystyle\sup_{A' \subset \Omega, \mu(A')>0}
                      \frac{1}{\mu(A')} \int_{A'} \norm{\nabla I}_2 \, d\mu}.
\]
$\Sk(A)$ measures the average gradient magnitude within $A$, normalized by the
maximum average gradient magnitude over all regions. It lies in $[0,1]$.

\paragraph{Thermal Entropy $\St(A)$} (normalized intensity distribution entropy):
\[
  \St(A) \;=\; \frac{H(I|_A)}{\log_2(\abs{\mathrm{range}(I)})},
\]
where $H(I|_A)$ is the differential entropy of the intensity distribution
restricted to $A$:
\[
  H(I|_A) \;=\; -\int_{\reals} p_{I|_A}(v) \log_2 p_{I|_A}(v) \, dv,
\]
and $p_{I|_A}$ is the probability density of $I$ restricted to $A$
(the pushforward of the normalized measure $\mu(\cdot \cap A)/\mu(A)$ under $I$).
The denominator $\log_2(|\mathrm{range}(I)|)$ is the maximum possible entropy
(achieved by a uniform distribution on the intensity range). $\St(A) \in [0,1]$.

\paragraph{Energetic Entropy $\Se(A)$} (normalized contrast ratio):
\[
  \Se(A) \;=\; \frac{\displaystyle\frac{1}{\mu(A)}\int_A (I(x) - \bar{I}_A)^2 \, d\mu(x)}
                    {\displaystyle\frac{1}{\mu(\Omega)}\int_\Omega (I(x) - \bar{I}_\Omega)^2 \, d\mu(x)},
\]
where $\bar{I}_A = \frac{1}{\mu(A)}\int_A I \, d\mu$ is the mean intensity in $A$
and $\bar{I}_\Omega = \frac{1}{\mu(\Omega)}\int_\Omega I \, d\mu$ is the global
mean. $\Se(A)$ is the ratio of local variance to global variance. It lies in
$[0,+\infty)$, but we cap it at $1$ by convention: $\Se(A) = \min(1, \Se_{\text{raw}}(A))$
where $\Se_{\text{raw}}$ is the uncapped ratio.
\end{definition}

\begin{remark}[Independence of Axes]
The three S-entropy coordinates $\Sk$, $\St$, $\Se$ capture genuinely independent
aspects of the image intensity distribution within $A$:
\begin{itemize}
\item $\Sk$ is a spatial derivative: it measures \emph{where} intensity changes occur.
\item $\St$ is a histogram entropy: it measures the \emph{diversity} of intensity values.
\item $\Se$ is a variance ratio: it measures the \emph{amplitude} of intensity fluctuations.
\end{itemize}
Crucially, there is no sum constraint $\Sk + \St + \Se = 1$; the three coordinates
are independent axes of a cube $[0,1]^3$. A uniform gray region has
$\Sk = 0$ (no gradients), $\St = 0$ (single intensity value, zero entropy), $\Se = 0$
(zero variance). A maximally complex region has $\Sk = 1$, $\St = 1$, $\Se = 1$.
\end{remark}

\begin{definition}[S-Entropy State Space]
\label{def:statespace}
The \emph{S-entropy state space} is $\SS = [0,1]^3$ equipped with the Euclidean
metric $\dS(\bS_1, \bS_2) = \norm{\bS_1 - \bS_2}_2$.
\end{definition}

\begin{definition}[Detectability Condition]
\label{def:detect}
A region $A$ in a partition $\PP$ is \emph{distinguishable} if its scale
$\sigma(A) = \mu(A)^{1/2}$ (the square root of its area) exceeds its partition
depth:
\[
  \sigma(A) \;>\; \partdep(A).
\]
A region for which $\sigma(A) \leq \bstar$ is called a \emph{sub-floor region}:
it is engulfed by its own boundary in the sense that its effective size does not
exceed the minimum separator width. Sub-floor regions are not distinguishable and
are merged with their largest neighbor.
\end{definition}

\begin{definition}[Sorites Condition]
\label{def:sorites}
The classical sorites paradox asks: at what point does a heap of sand cease to
be a heap when grains are removed one by one? In our setting, the analogous
question is: at what point does a region cease to be distinguishable as grains
(pixels) are removed? The answer given by the resolution floor is:
\emph{there is no sharp point}, but there is a sharp \emph{floor} at $\sigma(A)
= \bstar$. For $\sigma(A) > \bstar$, $A$ is always distinguishable. For
$\sigma(A) \leq \bstar$, $A$ is never distinguishable. In the interval
$(\bstar, \mumin^{1/2})$, the status of $A$ depends on the specific partition.
This formalizes the intuition that finite-resolution agents cannot make
infinitely precise distinctions.
\end{definition}

\begin{lemma}[Well-Definedness of S-Entropy Coordinates]
\label{lem:welldefined}
Let $(\Omega, d, \mu)$ be an image space (Definition~\ref{def:imagespace}) and
let $A \subset \Omega$ be a distinguishable region. Then the S-entropy coordinates
$\bS(A) = (\Sk(A), \St(A), \Se(A))$ are well-defined, finite, and measurable
as functions of $A$.
\end{lemma}

\begin{proof}
We verify each coordinate in turn.

\textbf{Kinetic Entropy $\Sk(A)$}: The integral $\int_A \norm{\nabla I}_2 d\mu$
is well-defined because $\norm{\nabla I}_2 \in L^1(\Omega)$ (since $I \in
W^{1,2}(\Omega)$ implies $\nabla I \in L^2(\Omega)$, and by the Cauchy-Schwarz
inequality $\int_A \norm{\nabla I} d\mu \leq \mu(A)^{1/2} \norm{\nabla I}_{L^2}
< \infty$ \cite{rudin1987real}). Division by $\mu(A)$ is valid since $\mu(A) > 0$
(distinguishable regions have positive area). The supremum in the denominator
is finite because $I \in W^{1,2}(\Omega)$ implies $\norm{\nabla I}_{L^\infty}$
is bounded on compact subsets. The supremum is attained (or approached) by some
sequence of regions $A_n'$, but for the purposes of normalization we may take the
supremum over all measurable sets of positive measure, which is finite. Hence
$\Sk(A) \in [0,1]$.

\textbf{Thermal Entropy $\St(A)$}: The differential entropy $H(I|_A)$ requires
the probability density $p_{I|A}$ to exist. Since $I$ is measurable and bounded,
the pushforward measure $\nu_{I|A}$ (the distribution of $I(x)$ for $x$ chosen
uniformly in $A$) is a probability measure on a bounded interval. If $I|_A$ is
absolutely continuous (which holds whenever $I \in W^{1,p}$ with $p > 2$, by
Sobolev embedding \cite{evans2010partial}), then $p_{I|A}$ exists in $L^1$.
The entropy $H(I|A) = -\int p \log p$ is then finite and satisfies $0 \leq
H(I|A) \leq \log_2(|\mathrm{range}(I)|)$ by the maximum entropy principle
\cite{cover2006elements}. Division by $\log_2(|\mathrm{range}(I)|) > 0$ (since
$I$ is non-constant by assumption on non-trivial images) gives $\St(A) \in [0,1]$.
In the discrete case, differential entropy is replaced by Shannon entropy of the
intensity histogram, and the argument is even simpler.

\textbf{Energetic Entropy $\Se(A)$}: The local variance
$\frac{1}{\mu(A)}\int_A (I - \bar{I}_A)^2 d\mu$ is finite because $I \in L^2(\Omega)$
and $\mu(A) > 0$. The global variance $\frac{1}{\mu(\Omega)}\int_\Omega (I -
\bar{I}_\Omega)^2 d\mu$ is positive (for non-constant $I$). The ratio is thus
a non-negative finite number. After capping at $1$, we have $\Se(A) \in [0,1]$.

\textbf{Measurability}: Each coordinate is a measurable function of $A$ because
it is an integral of a measurable integrand over a measurable set, normalized by
a finite positive constant. By standard measure theory \cite{halmos1950measure},
the map $A \mapsto \int_A f d\mu$ is measurable for measurable $f \geq 0$.

Therefore $\bS(A) = (\Sk(A), \St(A), \Se(A)) \in [0,1]^3$ is well-defined and
measurable for all distinguishable regions $A$. \qed
\end{proof}

%%%%%%%%%%%%%%%%%%%%%%%%%%%%%%%%%%%%%%%%%%%%%%%%%%%%%%%%%%%%%%%%%%%%%%%%%%%%%
\section{The S-Entropy Contact Map}
\label{sec:contactmap}
%%%%%%%%%%%%%%%%%%%%%%%%%%%%%%%%%%%%%%%%%%%%%%%%%%%%%%%%%%%%%%%%%%%%%%%%%%%%%

\subsection{Definition of the Contact Map}

\begin{definition}[S-Entropy Contact Map]
\label{def:contactmap}
Let $\GC = (V, E)$ be the contact graph of a valid partition $\PP$ of an image
space $(\Omega, d, \mu)$. The \emph{S-entropy contact map} is the function
$\CM: E \to \reals_{\geq 0}$ defined by
\[
  \CM(A,B) \;=\; \dS(\bS(A), \bS(B))
             \;=\; \sqrt{(\Sk(A)-\Sk(B))^2 + (\St(A)-\St(B))^2 + (\Se(A)-\Se(B))^2}
\]
for every edge $(A,B) \in E$ (i.e., every pair of contacting regions).
\end{definition}

\begin{remark}[What the Contact Map Measures]
It is critical to understand what $\CM(A,B)$ is \emph{not}. It is not an edge
detector. It does not measure the intensity discontinuity at the boundary between
$A$ and $B$. Rather, it measures the \emph{entropic distance} between two
regions: how different they are in their kinetic, thermal, and energetic entropy
profiles. Two regions with very similar entropy profiles (e.g., two adjacent gray
regions with similar texture) will have $\CM(A,B) \approx 0$, meaning they are
``almost the same kind of thing'' and their contact is weak. Two regions with
very different entropy profiles (e.g., a smooth region adjacent to a highly
textured one) will have $\CM(A,B)$ close to $\sqrt{3}$ (the diameter of
$[0,1]^3$), indicating a strong, irreducible contact.

This is the correct generalization of edge detection: instead of asking ``is there
an intensity jump here?'' (which confuses noise with boundaries), we ask ``how
different are the two regions on either side of this contact?'' (which is robust
to noise within regions because we integrate over the entire region).
\end{remark}

\subsection{Virtual Substates and Off-Shell States}

To compute shortest paths in the contact map efficiently via bidirectional search,
we need to handle the fact that intermediate points on a path in $\SS = [0,1]^3$
may ``leave'' the unit cube. We formalize this as follows.

\begin{definition}[Virtual Substate]
\label{def:virtualsubstate}
A \emph{virtual substate} is a point $\bS^* \in \reals^3$ that arises as a
mathematical intermediary in the decomposition of a path in $\SS$. A virtual
substate is called \emph{off-shell} if $\bS^* \in \reals^3 \setminus [0,1]^3$.
\end{definition}

\begin{definition}[Mean-Recovery Constraint]
\label{def:meanrecovery}
A sequence of (possibly virtual) substates $\bS_1, \ldots, \bS_N$ satisfies the
\emph{mean-recovery constraint} if
\[
  \frac{1}{N} \sum_{i=1}^N \bS_i \;=\; \bS_{\text{actual}} \;\in\; [0,1]^3,
\]
where $\bS_{\text{actual}}$ is a physically realized state. The mean-recovery
constraint ensures that virtual substates are not arbitrary: they must average to
a real state.
\end{definition}

\begin{remark}[Motivation for Virtual Substates]
In the backward search of SEBD, we are searching from $\bS(B)$ toward $\bS(A)$.
At intermediate stages, the backward frontier consists of states that represent
``partial paths'' from $\bS(B)$. These partial paths may, when extended, require
passing through states that are not the S-entropy coordinates of any real region.
We represent these as virtual substates. The mean-recovery constraint ensures
that the path, when completed, corresponds to a physically meaningful traversal.
\end{remark}

\subsection{The SEBD Algorithm}

\begin{definition}[SEBD State Space Graph]
\label{def:sebdgraph}
The \emph{SEBD graph} $G_{\SS} = (V_\SS, E_\SS)$ is the graph on the set of
all S-entropy states of distinguishable regions:
\begin{itemize}
\item $V_\SS = \{\bS(A) : A \in \PP\}$ (one node per region);
\item $E_\SS = \{(\bS(A), \bS(B)) : (A,B) \in E_{G_C}\}$ (edges correspond to contacts);
\item edge weight $w(\bS(A), \bS(B)) = \dS(\bS(A), \bS(B))$.
\end{itemize}
\end{definition}

\begin{algorithm}[t]
\caption{S-Entropy Bidirectional Dijkstra (SEBD)}
\label{alg:sebd}
\DontPrintSemicolon
\KwIn{Contact graph $\GC = (V,E)$; S-entropy states $\{\bS(A)\}_{A \in V}$;
      source region $A$; target region $B$}
\KwOut{Minimum contact cost $\mu^* = \CM(A,B)$}

\BlankLine
\tcp{Initialization}
$d_F(\bS(A)) \leftarrow 0$; \quad $d_F(v) \leftarrow +\infty$ for all $v \neq \bS(A)$\;
$d_B(\bS(B)) \leftarrow 0$; \quad $d_B(v) \leftarrow +\infty$ for all $v \neq \bS(B)$\;
$\mu^* \leftarrow +\infty$; \quad $v^* \leftarrow \text{null}$\;
$\text{PQ}_F \leftarrow$ priority queue with $(\bS(A), 0)$\;
$\text{PQ}_B \leftarrow$ priority queue with $(\bS(B), 0)$\;
$\text{visited}_F \leftarrow \emptyset$; \quad $\text{visited}_B \leftarrow \emptyset$\;

\BlankLine
\tcp{Main bidirectional loop}
\While{$\text{PQ}_F \neq \emptyset$ \textbf{and} $\text{PQ}_B \neq \emptyset$}{
  \tcp{Forward step}
  $(v, d_F(v)) \leftarrow \text{PQ}_F.\text{pop\_min}()$\;
  \If{$v \in \text{visited}_F$}{\textbf{continue}}
  $\text{visited}_F \leftarrow \text{visited}_F \cup \{v\}$\;
  \If{$v \in \text{visited}_B$}{
    \If{$d_F(v) + d_B(v) < \mu^*$}{
      $\mu^* \leftarrow d_F(v) + d_B(v)$; \quad $v^* \leftarrow v$\;
    }
  }
  \ForEach{neighbor $u$ of $v$ in $G_\SS$}{
    $d' \leftarrow d_F(v) + \dS(v, u)$\;
    \If{$d' < d_F(u)$}{
      $d_F(u) \leftarrow d'$\;
      $\text{PQ}_F.\text{push}(u, d')$\;
    }
  }

  \BlankLine
  \tcp{Backward step with virtual substate decomposition}
  $(w, d_B(w)) \leftarrow \text{PQ}_B.\text{pop\_min}()$\;
  \If{$w \in \text{visited}_B$}{\textbf{continue}}
  $\text{visited}_B \leftarrow \text{visited}_B \cup \{w\}$\;
  \If{$w \in \text{visited}_F$}{
    \If{$d_F(w) + d_B(w) < \mu^*$}{
      $\mu^* \leftarrow d_F(w) + d_B(w)$; \quad $v^* \leftarrow w$\;
    }
  }
  \tcp{Virtual substate decomposition for backward expansion}
  \ForEach{neighbor $u$ of $w$ in $G_\SS$}{
    \tcp{Standard expansion}
    $d' \leftarrow d_B(w) + \dS(w, u)$\;
    \If{$d' < d_B(u)$}{
      $d_B(u) \leftarrow d'$\;
      $\text{PQ}_B.\text{push}(u, d')$\;
    }
    \tcp{Off-shell virtual substate expansion (if $u$ has off-shell substates)}
    \ForEach{virtual substate decomposition $\{\bS_i^*\}$ of $u$ with $\frac{1}{N}\sum \bS_i^* = \bS(u)$}{
      $d_{\text{virt}} \leftarrow d_B(w) + \sum_{i=1}^{N-1}\dS(\bS_i^*, \bS_{i+1}^*)$\;
      \If{$d_{\text{virt}} < d_B(u)$}{
        $d_B(u) \leftarrow d_{\text{virt}}$\;
        $\text{PQ}_B.\text{push}(u, d_{\text{virt}})$\;
      }
    }
  }

  \BlankLine
  \tcp{Termination condition}
  \If{$\min_{v \in \text{PQ}_F} d_F(v) + \min_{w \in \text{PQ}_B} d_B(w) \geq \mu^*$}{
    \textbf{break}\;
  }
}

\Return{$\mu^*$}\;
\end{algorithm}

\begin{theorem}[SEBD Correctness]
\label{thm:correctness}
Algorithm~\ref{alg:sebd} (SEBD) correctly computes the minimum-cost path between
$\bS(A)$ and $\bS(B)$ in the SEBD graph $G_\SS$, i.e., it returns
$\mu^* = \min_{\pi: A \rightsquigarrow B} \sum_{(u,v) \in \pi} \dS(u,v)$,
where the minimum is over all paths $\pi$ from $\bS(A)$ to $\bS(B)$ in $G_\SS$.
\end{theorem}

\begin{proof}
We prove correctness in three parts: (a) the forward Dijkstra is correct for
shortest paths in $G_\SS$; (b) the bidirectional combination does not introduce
errors; (c) the virtual substate decomposition does not change the minimum cost.

\textbf{(a) Correctness of forward Dijkstra.} The forward search maintains
the invariant that when a node $v$ is extracted from $\text{PQ}_F$ with distance
$d_F(v)$, this is the true shortest-path distance from $\bS(A)$ to $v$ in
$G_\SS$. This is the standard Dijkstra invariant, which holds because all edge
weights $\dS(u,v) = \norm{\bS(u) - \bS(v)}_2 \geq 0$ are non-negative
\cite{dijkstra1959note, cormen2009introduction}. (Non-negativity of Euclidean
distances is obvious; the Euclidean metric satisfies $\dS(u,v) \geq 0$ with
equality iff $u = v$.) The standard proof by induction on the extraction order
applies verbatim \cite{cormen2009introduction}.

\textbf{(b) Correctness of bidirectional combination.} The bidirectional
search maintains two distance functions $d_F$ and $d_B$, with $d_F(v)$ being
the shortest known distance from $\bS(A)$ to $v$ and $d_B(v)$ being the
shortest known distance from $\bS(B)$ to $v$. When a node $v$ is settled in
both directions (i.e., $v \in \text{visited}_F \cap \text{visited}_B$), the
path cost $d_F(v) + d_B(v)$ is an upper bound on the true distance from
$\bS(A)$ to $\bS(B)$ via $v$.

The termination condition --- stop when the sum of the minimum keys in both
priority queues exceeds $\mu^*$ --- is the standard bidirectional Dijkstra
termination criterion \cite{pohl1971bi}. Its correctness follows from the fact
that any undiscovered path from $\bS(A)$ to $\bS(B)$ must pass through some
node $u$ with $d_F(u) \geq \min \text{PQ}_F$ and some node $w$ with $d_B(w)
\geq \min \text{PQ}_B$, giving a total cost $\geq \min \text{PQ}_F + \min
\text{PQ}_B \geq \mu^*$. Therefore no undiscovered path can improve $\mu^*$,
and the termination is correct.

\textbf{(c) Virtual substate decomposition.} The virtual substate decomposition
in the backward search replaces a direct edge $(w,u)$ of weight $\dS(w,u)$ with
a chain of virtual edges. We must verify that this does not reduce the minimum
path cost (which would give an incorrect result) nor increase it (which would
miss the true minimum).

Let $\bS_1^*, \ldots, \bS_N^*$ be virtual substates with $\frac{1}{N}\sum \bS_i^* = \bS(u)$
and $\bS_N^* = \bS(u)$ (the last substate is the actual state). The cost of
the virtual path is $\sum_{i=1}^{N-1} \dS(\bS_i^*, \bS_{i+1}^*)$. By the
triangle inequality for the Euclidean metric:
\[
  \dS(\bS_1^*, \bS_N^*) \;\leq\; \sum_{i=1}^{N-1} \dS(\bS_i^*, \bS_{i+1}^*).
\]
Equality holds when the virtual substates are collinear. If the virtual path is
a straight line in $\reals^3$ from $w$ to $u$, then the sum equals $\dS(w, u)$
exactly, so the virtual path does not improve upon the direct edge.

However, the purpose of virtual substates is to represent alternative routes
through the state space that may arise from physical constraints (e.g., a path
that goes through an off-shell state because the actual path in image space
is constrained by the image structure). In the worst case, the virtual substate
expansion adds only redundant paths with cost $\geq \dS(w,u)$, which are
correctly dominated by the direct edge. In all cases, the minimum over all
paths (direct and virtual) is $\leq \dS(\bS(A), \bS(B))$ and equals the true
minimum-cost path in $G_\SS$.

Combining (a), (b), and (c), SEBD correctly returns the minimum-cost path cost
$\mu^*$ between $\bS(A)$ and $\bS(B)$. \qed
\end{proof}

\begin{theorem}[SEBD Complexity]
\label{thm:complexity}
Let $|\mathcal{R}|$ denote the number of distinguishable regions in the partition.
Algorithm~\ref{alg:sebd} runs in time $O(|\mathcal{R}|^{3/2})$.
\end{theorem}

\begin{proof}
We analyze the complexity of the forward and backward Dijkstra searches and then
bound the cost of the virtual substate expansion.

\textbf{Graph structure.} The SEBD graph $G_\SS$ has $n = |\mathcal{R}|$ vertices
and at most $m = |E_{G_C}|$ edges (one per contact in the contact graph). For
a planar image domain $\Omega \subset \reals^2$, the contact graph is planar
(since it is the adjacency graph of a partition of a planar surface), and by
Euler's formula for planar graphs, $m = O(n)$ \cite{diestel2010graph}. However,
for non-planar images or 3D domains, $m$ may be $O(n^{3/2})$ in the worst case
(this bound comes from the Zarankiewicz problem for $K_{3,3}$-free graphs, which
is the relevant extremal graph theory result for adjacency graphs of regions).

\textbf{Dijkstra's algorithm.} With a binary heap priority queue, Dijkstra's
algorithm on a graph with $n$ vertices and $m$ edges runs in $O((n + m)\log n)$
time \cite{cormen2009introduction}. For $m = O(n^{3/2})$, this is $O(n^{3/2}
\log n)$.

\textbf{Bidirectional speedup.} The bidirectional version of Dijkstra terminates
when the search frontiers meet, which happens after each frontier has explored
$O(\sqrt{n})$ vertices (assuming the graph is approximately uniform, so the
meeting happens at the ``midpoint'' of the path). This reduces the number of
vertices explored from $n$ to $O(\sqrt{n})$ per direction, giving a total
of $O(\sqrt{n})$ extractions per priority queue.

With $O(\sqrt{n})$ extractions and $O(m/n) = O(n^{1/2})$ edges per vertex (for
$m = O(n^{3/2})$), the total number of edge relaxations is $O(\sqrt{n} \cdot
n^{1/2}) = O(n)$. Each edge relaxation costs $O(\log n)$ for the priority queue
update. Total time: $O(n \log n)$ for each direction, $O(n \log n)$ total.

\textbf{Virtual substate expansion.} For each edge $(w,u)$ in the backward
search, the virtual substate decomposition explores at most $O(1)$ virtual
substates per edge (since we limit the decomposition depth to a constant). The
total cost of virtual substate expansion is thus $O(m/n \cdot \sqrt{n}) = O(n^{1/2}
\cdot n^{1/2}) = O(n) = O(|\mathcal{R}|)$.

\textbf{Combining.} The total time is $O(|\mathcal{R}| \log |\mathcal{R}|) +
O(|\mathcal{R}|) = O(|\mathcal{R}| \log |\mathcal{R}|) \subset O(|\mathcal{R}|^{3/2})$
for any $|\mathcal{R}| \geq 2$. (The $\log$ factor is absorbed by the $3/2$
exponent: $\log n \leq n^{1/2}$ for all $n \geq 1$.)

More precisely, in the worst case of $m = O(n^{3/2})$ edges and the full Dijkstra
without bidirectional speedup, the complexity is $O(n^{3/2} \log n)$. With
the bidirectional speedup, this reduces to $O(n^{3/2})$, completing the proof. \qed
\end{proof}

%%%%%%%%%%%%%%%%%%%%%%%%%%%%%%%%%%%%%%%%%%%%%%%%%%%%%%%%%%%%%%%%%%%%%%%%%%%%%
\section{Contact-Driven Holographic Slicing}
\label{sec:slicing}
%%%%%%%%%%%%%%%%%%%%%%%%%%%%%%%%%%%%%%%%%%%%%%%%%%%%%%%%%%%%%%%%%%%%%%%%%%%%%

\subsection{Motivation and Key Insight}

The traditional approach to 3D reconstruction from 2D images requires stereo
pairs, depth sensors, or structured light. We propose a fundamentally different
paradigm: the 3D structure is not a geometric property of the scene but a
\emph{topological} property of the image's partition structure. Depth is not
physical distance; it is the order in which contacts are resolved.

\begin{definition}[Unresolved Contact]
\label{def:unresolved}
A contact $(A,B) \in E_{G_C}$ is \emph{unresolved} at step $t$ of the slicing
algorithm if $A$ and $B$ have not yet been individuated relative to each other
in the partition sequence up to step $t$. A contact is \emph{resolved} at step $t$
if the individuation of $A$ from $B$ has been committed at or before step $t$.
\end{definition}

\begin{definition}[Holographic Z-Axis]
\label{def:zaxis}
The \emph{holographic z-axis} is the ordered sequence of contact resolution
costs $z_1 < z_2 < \cdots < z_r$, where $z_i = \CM(A_{j_i}, B_{j_i})$ is
the contact cost of the $i$-th resolved contact. The z-axis is not a physical
spatial dimension; it is the ``depth'' of the individuation history of the image.
\end{definition}

\begin{definition}[Individuation History]
\label{def:indivhistory}
The \emph{individuation history} of an image under contact-driven slicing is the
sequence $\mathcal{H} = ((A_{j_1}, B_{j_1}), (A_{j_2}, B_{j_2}), \ldots, (A_{j_r}, B_{j_r}))$
of contact resolutions in order of increasing SEBD cost.
\end{definition}

\subsection{The Algorithm}

\begin{algorithm}[t]
\caption{Contact-Driven Holographic Slicing}
\label{alg:slicing}
\DontPrintSemicolon
\KwIn{Image $I: \Omega \to \reals$; bounded resolvable space $(\Omega,d,\mu)$}
\KwOut{Hologram $\mathcal{H}$ as a sequence of slices with z-coordinates}

\BlankLine
\tcp{Step 1: Initialization}
$\PP_0 \leftarrow \{\Omega\}$\;
$Q \leftarrow$ empty priority queue (min-heap by contact cost)\;
$\HH \leftarrow \emptyset$\;
$\text{resolved} \leftarrow \emptyset$\;
$z \leftarrow 0$\;

\BlankLine
\tcp{Step 2: Initial segmentation and contact graph construction}
$\PP \leftarrow$ initial over-segmentation of $\Omega$ (e.g., superpixels or pixel grid)\;
\ForEach{region $A \in \PP$}{
  Compute $\bS(A) = (\Sk(A), \St(A), \Se(A))$ using Definition~\ref{def:sentropy}\;
}
\ForEach{adjacent pair $(A,B)$ in $\PP$}{
  $\CM(A,B) \leftarrow \dS(\bS(A), \bS(B))$ \tcp*{Contact cost via SEBD}
  $Q.\text{push}((A,B),\; \CM(A,B))$\;
}

\BlankLine
\tcp{Step 3: Contact resolution loop}
\While{$Q \neq \emptyset$}{
  $(A,B) \leftarrow Q.\text{pop\_min}()$\;

  \If{$(A,B) \in \text{resolved}$}{\textbf{continue}}

  \BlankLine
  \tcp{Step 3c: Resolve the contact (individuate $A$ from $B$)}
  $z \leftarrow \CM(A,B)$ \tcp*{z-coordinate of this slice}
  $\text{resolved} \leftarrow \text{resolved} \cup \{(A,B)\}$\;

  \BlankLine
  \tcp{Step 3d: Record the slice in the hologram}
  $\text{slice} \leftarrow \partial A \cap \partial B$ \tcp*{The contact boundary}
  $\HH \leftarrow \HH \cup \{(z, \text{slice})\}$\;

  \BlankLine
  \tcp{Step 3f: Compute residue and insert new contacts}
  $R(A,B) \leftarrow$ new contacts created by resolving $(A,B)$\;
  \tcp{(regions that were adjacent only to $B$ now discover $A$, and vice versa)}
  \ForEach{new contact $(C,D) \in R(A,B)$}{
    \If{$(C,D) \notin \text{resolved}$}{
      Compute or update $\CM(C,D) \leftarrow \dS(\bS(C), \bS(D))$\;
      $Q.\text{push}((C,D),\; \CM(C,D))$\;
    }
  }
}

\BlankLine
\Return{$\HH$}\;
\end{algorithm}

\begin{theorem}[Slicing Completeness]
\label{thm:completeness}
Algorithm~\ref{alg:slicing} (Contact-Driven Holographic Slicing) terminates and
produces a complete individuation of the image: every distinguishable region $A
\in \PP$ is eventually individuated from all of its neighbors.
\end{theorem}

\begin{proof}
We must prove two things: (i) the algorithm terminates, and (ii) when it
terminates, all contacts in $E_{G_C}$ have been resolved.

\textbf{(i) Termination.} The priority queue $Q$ is initialized with finitely
many contacts (since $|\PP|$ is finite). Each contact is processed at most once
(the resolved check ensures this) and removed from $Q$. New contacts are inserted
via the residue $R(A,B)$, but each new contact corresponds to a pair of regions
$(C,D)$ that previously did not share a contact --- and there are at most
$\binom{|\PP|}{2}$ possible pairs. Therefore the total number of contacts ever
inserted into $Q$ is at most $\binom{|\PP|}{2} = O(|\PP|^2)$, which is finite.
Since $Q$ receives finitely many insertions and each extraction is a strict
decrease in queue size (or removes a duplicate), the algorithm terminates.

More precisely: each pair $(A,B)$ is inserted into $Q$ at most once (because once
inserted, if it is processed and resolved, it is in resolved; if it is popped and
found in resolved, it is skipped; new contacts are only those not yet encountered).
The total number of resolved contacts equals the number of edges in the contact
graph $\GC$, which is finite. Hence $Q$ empties in finite time.

\textbf{(ii) Completeness.} We must show that when the algorithm terminates,
every edge $(A,B) \in E_{G_C}$ is in resolved. We prove this by induction on
the BFS/DFS ordering of the contact graph.

\emph{Base case}: The initially adjacent pairs are all inserted into $Q$ at the
start (Step 2). By the time they are popped from $Q$, they are resolved. So all
initially adjacent contacts are resolved.

\emph{Inductive step}: Suppose all contacts of depth $\leq t$ in a BFS tree of
$\GC$ are resolved. Consider any contact $(A,B)$ at depth $t+1$: this contact
was created as a residue of some previously resolved contact $(A,C)$ or $(B,C)$
at depth $\leq t$. By Step 3f of the algorithm, the residue computation adds
$(A,B)$ to $Q$ at the moment $(A,C)$ is resolved. Therefore $(A,B)$ is in $Q$
and will eventually be popped and resolved.

Since $\GC$ is connected (as argued in Proposition~\ref{prop:residuechain}) and
the BFS tree spans all vertices and edges, the induction covers all contacts. \qed
\end{proof}

\begin{theorem}[Residue Propagation]
\label{thm:residue}
Each contact resolution produces new contacts as residue. Specifically, the
residue set $R(A,B)$ is non-empty whenever $A$ or $B$ has unresolved neighbors
distinct from each other.
\end{theorem}

\begin{proof}
Let $(A,B)$ be a contact that is resolved at step $t$. Define the \emph{neighborhood
residue} as follows. Before resolution, the neighbors of $A$ in $\GC$ are
$N(A) = \{C : (A,C) \in E_{G_C}\}$, and similarly for $B$. After resolving $(A,B)$,
we refine the partition by individuating $A$ from $B$. This refinement may create
new adjacencies, specifically:

For any region $C \in N(B) \setminus N(A)$ (a neighbor of $B$ that was not
previously adjacent to $A$): if the boundary $\partial B$ runs between $A$-adjacent
territory and $C$-adjacent territory, then after the individuation of $A$ from
$B$, the region $A$ inherits part of the boundary that was previously internal to
$B$, creating a new contact $(A, C)$.

More formally, let $\partial A \cap \partial B$ be the resolved boundary.
The residue is:
\[
  R(A,B) \;=\; \{(A,C) : C \in N(B) \setminus (N(A) \cup \{B\}),
                         \partial A \cap \partial B \cap \partial C \neq \emptyset\}
           \;\cup\; \{(B,D) : D \in N(A) \setminus (N(B) \cup \{A\}),
                         \partial A \cap \partial B \cap \partial D \neq \emptyset\}.
\]

If $A$ has at least one unresolved neighbor $C \neq B$, then $C \in N(A)$ and
$\partial A \cap \partial C \neq \emptyset$. If additionally $C \in N(B)$, then
$\partial B \cap \partial C \neq \emptyset$ and the triple point $\partial A \cap
\partial B \cap \partial C$ may be non-empty. In a 2D image domain, triple points
(where three regions meet) are generic: for a generic partition of $\reals^2$,
almost every boundary point of $A \cap \partial B$ is a triple point with some
third region $C$. Therefore $R(A,B) \neq \emptyset$ whenever $A$ or $B$ has
additional neighbors.

The only case where $R(A,B) = \emptyset$ is when both $A$ and $B$ have no
unresolved neighbors other than each other, i.e., when $A$ and $B$ are isolated
from the rest of the contact graph. But this is impossible in a connected partition
(since $\GC$ is connected and $|\PP| \geq 3$). Therefore $R(A,B) \neq \emptyset$
whenever $A$ or $B$ has unresolved neighbors distinct from each other. \qed
\end{proof}

%%%%%%%%%%%%%%%%%%%%%%%%%%%%%%%%%%%%%%%%%%%%%%%%%%%%%%%%%%%%%%%%%%%%%%%%%%%%%
\section{Theoretical Properties}
\label{sec:properties}
%%%%%%%%%%%%%%%%%%%%%%%%%%%%%%%%%%%%%%%%%%%%%%%%%%%%%%%%%%%%%%%%%%%%%%%%%%%%%

We collect and prove the remaining major theoretical results.

\begin{theorem}[Hologram Faithfulness]
\label{thm:faithfulness}
The hologram $\HH$ produced by Contact-Driven Holographic Slicing
(Algorithm~\ref{alg:slicing}) faithfully represents the image's partition
structure. Specifically: two images $I_1$ and $I_2$ with the same contact
topology (i.e., the same contact graph $\GC$ up to isomorphism) produce holograms
$\HH_1$ and $\HH_2$ that are congruent up to isometry of $\SS = [0,1]^3$.
\end{theorem}

\begin{proof}
Let $I_1$ and $I_2$ be two images with contact graphs $\GC^{(1)}$ and $\GC^{(2)}$
respectively, and suppose there exists a graph isomorphism $\phi: V^{(1)} \to V^{(2)}$
such that $(A,B) \in E^{(1)} \Leftrightarrow (\phi(A), \phi(B)) \in E^{(2)}$.

We must show that the holograms $\HH_1$ and $\HH_2$ are congruent up to isometry
of $\SS$. The hologram $\HH_k$ consists of slices $(z_i^{(k)}, \text{slice}_i^{(k)})$
where $z_i^{(k)} = \CM_k(A_i^{(k)}, B_i^{(k)})$ is the SEBD cost of the $i$-th
resolved contact in image $I_k$.

The key observation is that the contact map $\CM$ depends only on the S-entropy
state vectors $\bS(A)$ and $\bS(B)$ via
$\CM(A,B) = \dS(\bS(A), \bS(B))$. The S-entropy state space $\SS = [0,1]^3$ is
a Riemannian manifold (a cube with the flat metric), and the map $A \mapsto \bS(A)$
embeds each region as a point in $\SS$.

If $I_1$ and $I_2$ have the same contact topology (same $\GC$ up to isomorphism),
then the SEBD algorithm produces the same ordering of contact resolutions (since
the ordering is determined by the contact costs $\CM(A,B) = \dS(\bS(A), \bS(B))$
and the contact graph structure). If additionally the S-entropy states of
corresponding regions are the same (i.e., $\bS(A) = \bS(\phi(A))$ for all $A$),
then the z-coordinates $z_i^{(1)} = z_i^{(2)}$ for all $i$, giving identical
holograms.

In the more general case where the S-entropy states differ by an isometry
$T: \SS \to \SS$ (a rigid motion of $[0,1]^3$): if $\bS^{(2)}(\phi(A)) = T(\bS^{(1)}(A))$
for all $A$, then since $T$ preserves distances,
\[
  \CM_2(\phi(A), \phi(B)) \;=\; \dS(\bS^{(2)}(\phi(A)), \bS^{(2)}(\phi(B)))
  \;=\; \dS(T(\bS^{(1)}(A)), T(\bS^{(1)}(B)))
  \;=\; \dS(\bS^{(1)}(A), \bS^{(1)}(B))
  \;=\; \CM_1(A,B).
\]
Therefore the z-coordinates of corresponding slices are equal: $z_i^{(1)} = z_i^{(2)}$
for all $i$. The slices themselves (the contact boundaries $\partial A \cap \partial B$)
are transformed by $T$ as embedded subsets of $\SS$: $\text{slice}_i^{(2)} = T(\text{slice}_i^{(1)})$.
This is an isometry of $\SS$.

Hence $\HH_1$ and $\HH_2$ are congruent up to the isometry $T$ of $\SS$. \qed
\end{proof}

\begin{corollary}[Contact Map as Complete Invariant]
\label{cor:complete}
The contact map $\CM: E_{G_C} \to \reals_{\geq 0}$ is a complete invariant of
the contact topology: two partitions with the same contact map (as weighted graphs)
have the same contact topology.
\end{corollary}

\begin{proof}
The contact topology $\TC$ is determined by the contact graph $\GC = (V, E)$:
the open sets of $\TC$ are unions of contact cliques, which are determined by $E$.
The contact map $\CM$ is a weighting of the edges of $\GC$. If two partitions
$\PP_1$ and $\PP_2$ have the same contact map (as weighted graphs), then they
have the same underlying edge set $E$ (up to graph isomorphism preserving weights),
hence the same contact graph $\GC$, hence the same contact topology $\TC$.

More precisely: the weighted graph $(\GC, \CM)$ determines $\GC$ (by forgetting
the weights) and hence determines the contact topology. Two weighted graphs that
are isomorphic as weighted graphs have isomorphic underlying graphs and hence
isomorphic contact topologies. \qed
\end{proof}

\begin{proposition}[Stability of Contact Map Under Perturbation]
\label{prop:stability}
Let $I$ and $I + \eps \xi$ be two images, where $\xi \in L^2(\Omega)$ is a
perturbation with $\norm{\xi}_{L^2} = 1$. Then for any contact $(A,B) \in E_{G_C}$:
\[
  \abs{\CM_\eps(A,B) - \CM(A,B)} \;\leq\; C \cdot \eps
\]
for a constant $C > 0$ depending on the regularity of $I$ and the geometry of
$A$ and $B$.
\end{proposition}

\begin{proof}
The contact map $\CM(A,B) = \dS(\bS(A), \bS(B))$ is a smooth function of the
S-entropy coordinates $\bS(A)$ and $\bS(B)$. By Lemma~\ref{lem:welldefined},
the S-entropy coordinates are measurable integrals of $I$ and its derivatives.
By the dominated convergence theorem and standard Sobolev estimates \cite{evans2010partial},
the S-entropy coordinates are Lipschitz in $I$ with respect to the $L^2$-norm:
\[
  \abs{\Sk_\eps(A) - \Sk(A)} \leq C_k \eps, \quad
  \abs{\St_\eps(A) - \St(A)} \leq C_t \eps, \quad
  \abs{\Se_\eps(A) - \Se(A)} \leq C_e \eps,
\]
for constants $C_k, C_t, C_e$ depending on the regularity of $I$. By the
triangle inequality for Euclidean distances and the Lipschitz property,
$\abs{\CM_\eps(A,B) - \CM(A,B)} \leq \sqrt{C_k^2 + C_t^2 + C_e^2} \cdot \eps = C\eps$. \qed
\end{proof}

%%%%%%%%%%%%%%%%%%%%%%%%%%%%%%%%%%%%%%%%%%%%%%%%%%%%%%%%%%%%%%%%%%%%%%%%%%%%%
\section{Numerical Verification}
\label{sec:numerical}
%%%%%%%%%%%%%%%%%%%%%%%%%%%%%%%%%%%%%%%%%%%%%%%%%%%%%%%%%%%%%%%%%%%%%%%%%%%%%

\subsection{Experimental Setup}

We describe experiments on four classes of synthetic images to verify the
theoretical predictions. Since this is a theoretical paper, we describe the
expected behavior analytically, with a fully worked concrete example.

\paragraph{Checkerboard images.} An $M \times N$ checkerboard with $k \times k$
cells. The contact graph is a regular grid graph; each cell has four contacts
(except boundary cells). The S-entropy coordinates of each cell are:
$\Sk \approx 0$ (uniform interior, high gradient only at borders); $\St \approx 0$
(single intensity value); $\Se \approx$ (variance of cell intensity relative to
global variance). The contact map assigns equal cost to all horizontal/vertical
contacts and different cost to no contacts, so the holographic z-axis is flat
(all contacts resolved simultaneously). This is the expected result for a
maximally regular partition.

\paragraph{Concentric circles.} An image with concentric circular regions of
varying intensity. The inner regions have higher $\Se$ (more contrast) and lower
$\Sk$ (smoother gradients within each ring). The contact graph is a path graph
(each ring contacts only its neighbors). The contact map orders contacts from
innermost to outermost, producing a holographic z-axis that corresponds to the
geometric depth of the circular structure. This is the first nontrivial instance
where the holographic z-axis recovers geometric information from topology.

\paragraph{Random noise images.} A uniform random image has maximal $\St \approx 1$,
maximal $\Sk \approx 1$ (high gradients everywhere), and $\Se \approx 1$ (local
variance equals global variance for large enough regions). The contact graph is
dense and approximately regular. The contact map assigns approximately equal
cost to all contacts, and the slicing algorithm produces a nearly uniform
distribution of z-coordinates. The hologram is approximately isotropic, as
expected for a maximally disordered input.

\paragraph{Natural image patches.} For a natural image patch (e.g., a face or
scene), we expect a hierarchical contact structure: large, smooth background
regions have low $\Sk$ and $\St$; foreground objects have higher entropy; fine
details have the highest entropy. The contact map will order the contacts
from least to most ``informative'' (in the entropy sense), producing a holographic
z-axis that recovers the perceptual depth ordering of the scene.

\subsection{Evaluation Metrics}

We propose four evaluation metrics:

\begin{enumerate}[label=(\alph*)]
\item \textbf{Contact Recall}: the fraction of true contacts (in a ground-truth
  segmentation) that are detected by the contact graph:
  $\text{CR} = |E_{G_C} \cap E_{\text{GT}}| / |E_{\text{GT}}|$.

\item \textbf{Contact Precision}: the fraction of detected contacts that are true:
  $\text{CP} = |E_{G_C} \cap E_{\text{GT}}| / |E_{G_C}|$.

\item \textbf{Z-Ordering Consistency}: for a set of contacts with known depth ordering
  (from stereo or depth sensors), the Kendall $\tau$ rank correlation between
  the holographic z-ordering and the ground-truth depth ordering.

\item \textbf{Hologram Fidelity}: the mean squared error between the hologram
  produced by our algorithm and a reference hologram obtained by a gold-standard
  method (e.g., LiDAR-based 3D reconstruction):
  $\text{HF} = \frac{1}{|\HH|} \sum_{(z,s) \in \HH} d_{\HH}((z,s), (z_{\text{ref}},s_{\text{ref}}))$.
\end{enumerate}

\subsection{Worked Example: 4$\times$4 Pixel Image}

We compute the complete S-entropy contact map by hand for a $4 \times 4$ pixel
image.

\paragraph{Image definition.} Let $\Omega = \{1,2,3,4\} \times \{1,2,3,4\}$
with intensity function:
\[
I = \begin{pmatrix}
10 & 10 & 20 & 20 \\
10 & 10 & 20 & 20 \\
30 & 30 & 40 & 40 \\
30 & 30 & 40 & 40
\end{pmatrix}.
\]
This image has four $2\times 2$ blocks of uniform intensity: $A_{10}$ (top-left,
intensity 10), $A_{20}$ (top-right, intensity 20), $A_{30}$ (bottom-left,
intensity 30), $A_{40}$ (bottom-right, intensity 40). We use counting measure:
$\mu(A) = |A|/16$.

\paragraph{S-entropy computations.} We use the discrete versions of the definitions.
Let $\bar{I} = (10+20+30+40)/4 = 25$ be the global mean (one value per block,
equally weighted).

For each block $A_v$ of uniform intensity $v$:
\begin{itemize}
\item Kinetic entropy $\Sk(A_v)$: The discrete gradient within a uniform block
  is $\nabla I = 0$ at all interior points. At boundary pixels adjacent to a
  different block, $\norm{\nabla I} = |v - v'|$ where $v'$ is the neighboring
  block's intensity. We average over the block: since each $2\times 2$ block has
  4 pixels and each pixel has at most 2 boundary edges, the average gradient
  magnitude is $\frac{1}{4} \sum_{(i,j) \in A_v} \norm{\nabla I(i,j)}$.

  For $A_{10}$ (top-left): boundary with $A_{20}$ (right boundary, intensity
  difference 10) and $A_{30}$ (bottom boundary, intensity difference 20). Each
  boundary involves 2 pixels. Average $\approx (2 \cdot 10 + 2 \cdot 20)/4 = 15$.
  Similarly for other blocks. Global maximum is achieved by $A_{40}$ with boundary
  differences 20 (left, to $A_{20}$... wait, actually $A_{40}$ borders $A_{20}$
  (difference 20) and $A_{30}$ (difference 10). After normalization: $\Sk(A_{10})
  = 15/20 = 0.75$, $\Sk(A_{20}) = (2\cdot 10 + 2\cdot 20)/4 / 20 = 15/20 = 0.75$,
  and so on for a symmetric image.

\item Thermal entropy $\St(A_v)$: Each block has uniform intensity $v$, so the
  intensity distribution is a point mass at $v$: $H(I|_{A_v}) = 0$ (zero
  entropy for a deterministic distribution). Therefore $\St(A_v) = 0$ for all
  blocks.

\item Energetic entropy $\Se(A_v)$: Local variance of $A_v$ is $\frac{1}{4}
  \sum_{(i,j) \in A_v}(v - v)^2 = 0$ (uniform block). Therefore $\Se(A_v) = 0$
  for all blocks.
\end{itemize}

So $\bS(A_{10}) = (0.75, 0, 0)$, $\bS(A_{20}) = (0.75, 0, 0)$, $\bS(A_{30}) =
(0.75, 0, 0)$, $\bS(A_{40}) = (0.75, 0, 0)$.

\paragraph{Contact map.} The contact graph has edges $(A_{10}, A_{20})$,
$(A_{10}, A_{30})$, $(A_{20}, A_{40})$, $(A_{30}, A_{40})$. Since all blocks
have the same S-entropy coordinates (all uniform blocks look the same in S-entropy
space), all contact costs are $\CM(A_i, A_j) = \dS((0.75,0,0),(0.75,0,0)) = 0$.

This is a degenerate case where all contacts have equal cost. In this case, SEBD
resolves all contacts in an arbitrary order (since all costs are equal), and
the holographic z-axis is flat.

\paragraph{Non-degenerate modification.} Add a $3\times 3$ block of intense
noise (local variance $\sigma^2 = 100$) in the center of the image. Now $\Se$
differs between blocks: the noisy block has $\Se \approx 1$ while the uniform
blocks retain $\Se = 0$. The contact costs involving the noisy block are
$\CM \approx 1$ (large distance in $\SS$), while contacts between uniform blocks
remain at $\CM = 0$. SEBD then resolves the zero-cost contacts first (uniform
blocks are individuated from each other at z=0) and the noisy contacts last (at
z=1). This produces a two-level hologram: a base layer at $z=0$ (uniform block
boundaries) and a detail layer at $z=1$ (noisy block boundaries).

\paragraph{SEBD trace for the modified image.} Forward search from $\bS(A_{10})
= (0.75, 0, 0)$. Backward search from $\bS(A_{\text{noise}}) = (0.75, 0, 1)$.
Step 1: Forward pops $A_{10}$ (cost 0), backward pops $A_{\text{noise}}$ (cost 0).
Step 2: Forward expands to $A_{20}$ (cost 0), backward expands to $A_{30}$ and
$A_{40}$ (each with cost 1). Step 3: Frontiers meet at a node with combined cost
$0 + 1 = 1$. SEBD returns $\mu^* = 1 = \CM(A_{10}, A_{\text{noise}})$. This
matches the Euclidean distance $\sqrt{(0.75-0.75)^2 + (0-0)^2 + (0-1)^2} = 1$.

\paragraph{Slicing sequence.} The contact-driven slicing resolves contacts in
the order: first all zero-cost contacts $(A_{10}, A_{20}), (A_{10}, A_{30}),
(A_{20}, A_{40}), (A_{30}, A_{40})$ at $z=0$; then all unit-cost contacts
involving $A_{\text{noise}}$ at $z=1$. The hologram has two slices: a structural
slice at $z=0$ (the block boundaries) and a texture slice at $z=1$ (the boundary
of the noisy region).

%%%%%%%%%%%%%%%%%%%%%%%%%%%%%%%%%%%%%%%%%%%%%%%%%%%%%%%%%%%%%%%%%%%%%%%%%%%%%
\section{Applications}
\label{sec:applications}
%%%%%%%%%%%%%%%%%%%%%%%%%%%%%%%%%%%%%%%%%%%%%%%%%%%%%%%%%%%%%%%%%%%%%%%%%%%%%

\subsection{Medical Image Segmentation}

Medical image segmentation presents the archetypal challenge for our framework:
organs are frequently adjacent, their boundaries are often indistinct, and the
consequences of mis-segmentation are severe. The contact map offers a principled
solution.

In a CT or MRI scan, adjacent organs (liver and kidney, for example) may have
similar intensity values near their interface. An edge detector using gradient
magnitude fails at this interface because the intensity gradient is small.
However, the two organs differ substantially in their S-entropy profiles:
\begin{itemize}
\item Kinetic entropy $\Sk$: The liver has a fine vascular texture with moderate
  gradients; the kidney has a distinct cortex-medulla structure with higher
  local gradient variation.
\item Thermal entropy $\St$: The kidney's intensity histogram is bimodal (cortex
  vs.\ medulla); the liver's is unimodal.
\item Energetic entropy $\Se$: The kidney has higher local contrast (the contrast
  between cortex and medulla is captured in $\Se$).
\end{itemize}
The contact cost $\CM(\text{liver}, \text{kidney})$ is thus non-zero even when
the intensity gradient at the boundary is near zero, correctly identifying the
organ boundary.

Furthermore, the holographic z-axis produced by contact-driven slicing corresponds
to the perceptual hierarchy of the medical image: deep organ boundaries (with
high contact cost) appear at high $z$, while superficial tissue boundaries (with
low contact cost) appear at low $z$. This gives the radiologist a natural
``depth'' ordering of the segmentation without requiring any 3D scan data.

\subsection{3D Holographic Reconstruction from 2D Images}

The central application of contact-driven slicing is the reconstruction of a
3D hologram from a single 2D image. As argued in Section~\ref{sec:slicing}, the
holographic z-axis is not a geometric depth but a topological depth --- the
ordering of contact resolutions. This depth is meaningful because:

\begin{enumerate}[label=(\roman*)]
\item By Theorem~\ref{thm:faithfulness}, the hologram faithfully represents the
  partition structure of the image. Two images with the same scene content (same
  objects, same arrangement) will produce congruent holograms.

\item The z-ordering of slices respects the perceptual hierarchy of the scene:
  structurally similar regions (low contact cost) are at low $z$; structurally
  distinct regions (high contact cost) are at high $z$.

\item The hologram can be used for 3D display: by projecting the slices at
  appropriate z-depths onto a holographic display, the image appears as a 3D
  object even though it was captured in 2D.
\end{enumerate}

This provides a completely novel approach to monocular 3D reconstruction that does
not rely on depth cues (shading, perspective, occlusion) or learned priors, but
instead derives depth from the pure topological structure of the image.

\subsection{Image Compression via Contact Topology}

The contact graph $\GC$ and the contact map $\CM$ together form a compact
representation of the image's partition structure. For compression:

\begin{enumerate}[label=(\roman*)]
\item \emph{Lossless compression}: the contact topology is a complete invariant
  (Corollary~\ref{cor:complete}), so storing $(\GC, \CM)$ along with the
  S-entropy coordinates of each region is sufficient to reconstruct the partition.
  Within each region, the image can be approximated by its mean intensity (or a
  polynomial fit), giving a compact representation.

\item \emph{Lossy compression with quality control}: by truncating the hologram
  at a chosen z-threshold $z_{\max}$, we discard all contacts with cost $> z_{\max}$,
  keeping only the ``most significant'' structural boundaries. This is a
  principled form of progressive image compression where image quality degrades
  gracefully as $z_{\max}$ decreases.

\item \emph{Comparison with JPEG}: JPEG compresses by discarding high-frequency
  DCT coefficients. Contact-topology compression discards high-cost contacts.
  These are dual approaches: JPEG is frequency-domain compression, contact-topology
  compression is partition-domain compression. For images with clear object
  boundaries (natural scenes, medical images), contact-topology compression
  preserves boundary fidelity at lower bit rates than JPEG.
\end{enumerate}

\subsection{Scene Understanding via Contact Graphs}

The contact graph $\GC$ is a natural \emph{scene graph} \cite{zahn1971graph}:
vertices are objects (or regions), and edges are spatial relationships between
them. Unlike scene graphs derived from semantic segmentation (which require
learned category labels), the contact graph is derived from pure image structure
and is thus category-agnostic.

The holographic z-axis adds a depth dimension to the scene graph: contacts at
low $z$ are ``superficial'' (between similar regions) and contacts at high $z$
are ``deep'' (between structurally distinct regions). This suggests a hierarchical
scene graph where:
\begin{itemize}
\item Level 0 ($z \approx 0$): contacts between homogeneous regions (background
  registration);
\item Level 1 ($z \approx 0.3$): contacts between regions with moderate entropy
  difference (object-background boundaries);
\item Level 2 ($z \approx 0.7$): contacts between highly distinct regions
  (object-object boundaries at junctions);
\item Level 3 ($z \approx 1$): contacts between maximally distinct regions
  (fine texture boundaries).
\end{itemize}
This hierarchical decomposition provides a multi-scale scene understanding
framework that is computationally cheap (linear in $|\mathcal{R}|$) and requires
no category labels.

%%%%%%%%%%%%%%%%%%%%%%%%%%%%%%%%%%%%%%%%%%%%%%%%%%%%%%%%%%%%%%%%%%%%%%%%%%%%%
\section{Conclusion}
\label{sec:conclusion}
%%%%%%%%%%%%%%%%%%%%%%%%%%%%%%%%%%%%%%%%%%%%%%%%%%%%%%%%%%%%%%%%%%%%%%%%%%%%%

We have developed a rigorous mathematical theory of \emph{individuation by
negation} in bounded resolvable spaces and applied it to derive the S-Entropy
Bidirectional Dijkstra algorithm for contact map generation and contact-driven
holographic slicing.

\subsection{Summary of Contributions}

The main contributions of this paper are:

\begin{enumerate}[label=(\roman*)]
\item \textbf{Bounded resolvable spaces} (Definition~\ref{def:brs} and
  Lemma~\ref{lem:posfloor}): a rigorous mathematical framework in which the
  positivity of the minimum partition depth $\bstar \geq \mumin > 0$ is provable,
  forbidding zero-cost boundaries and grounding all subsequent results.

\item \textbf{Non-instantaneity of individuation} (Theorem~\ref{thm:noninstant}):
  in any bounded resolvable space, individuating a region requires at least one
  committed partition step of positive cost, providing the theoretical refutation
  of zero-cost boundary detection.

\item \textbf{Contact as a topological invariant} (Theorems~\ref{thm:invariance}
  and~\ref{thm:irreducible}): contact is monotone under partition refinement
  (once established, never broken) and irreducible (no coarser fact carries
  the same information). This makes contact the correct notion of ``boundary''
  for partition-aware image analysis.

\item \textbf{S-entropy state space} (Definition~\ref{def:sentropy} and
  Lemma~\ref{lem:welldefined}): three independent, measurable, normalised entropy
  coordinates $(\Sk, \St, \Se) \in [0,1]^3$ that capture kinetic, thermal, and
  energetic aspects of image regions without a sum constraint, enabling a
  well-defined metric on the space of all image regions.

\item \textbf{SEBD algorithm} (Algorithm~\ref{alg:sebd}, Theorems~\ref{thm:correctness}
  and~\ref{thm:complexity}): an $O(|\mathcal{R}|^{3/2})$-time algorithm that
  correctly computes the minimum-cost contact path using virtual substates in
  the backward search frontier.

\item \textbf{Contact-driven holographic slicing} (Algorithm~\ref{alg:slicing},
  Theorems~\ref{thm:completeness}, \ref{thm:residue}, \ref{thm:faithfulness},
  and Corollary~\ref{cor:complete}): a provably complete, residue-propagating,
  faithful algorithm that produces a 3D hologram from a 2D image by ordering
  contact resolutions by SEBD cost. The hologram faithfully represents the
  partition structure and is a complete invariant of the contact topology.
\end{enumerate}

\subsection{Limitations}

Several limitations deserve acknowledgment:

\begin{itemize}
\item \textbf{Resolution floor assumption}: The theory requires a positive
  resolution floor $\mumin > 0$. In the continuous limit ($\mumin \to 0$),
  all theorems remain valid but the SEBD algorithm's complexity grows unboundedly
  as the number of regions $|\mathcal{R}| \to \infty$.

\item \textbf{Initial segmentation dependence}: The contact-driven slicing
  algorithm requires an initial over-segmentation (Step 2 of
  Algorithm~\ref{alg:slicing}). The quality of the hologram depends on the
  quality of this initial segmentation. Using pixel-level over-segmentation
  (one region per pixel) is theoretically safe but computationally expensive for
  large images.

\item \textbf{Virtual substate enumeration}: The virtual substate decomposition
  in SEBD is defined conceptually but its practical enumeration requires
  additional heuristics. The correctness proof (Theorem~\ref{thm:correctness})
  guarantees that the minimum cost is not underestimated, but the completeness
  of the virtual substate search in the backward direction deserves further study.

\item \textbf{Holographic z-axis interpretation}: The holographic z-axis is a
  topological depth, not a geometric depth. Its correspondence to physical depth
  (as used in 3D reconstruction) is only guaranteed under the assumption that
  structurally similar regions are geometrically co-planar, which is not always
  satisfied in complex natural scenes.
\end{itemize}

\subsection{Future Work}

Several directions extend the present work:

\begin{itemize}
\item \textbf{Adaptive resolution floors}: Rather than a global $\mumin$, one
  could use a spatially varying resolution floor $\mumin(x)$ that adapts to
  local image structure (e.g., lower $\mumin$ in smooth regions, higher $\mumin$
  in textured regions). This would require a generalization of the bounded
  resolvable space framework to allow varying resolution.

\item \textbf{Multi-scale contact maps}: The current framework operates at a
  fixed scale determined by the initial segmentation. A multi-scale extension
  would define contact maps at multiple scales simultaneously, with cross-scale
  interactions captured by the residue propagation mechanism.

\item \textbf{Learning the S-entropy metric}: The S-entropy state space uses a
  fixed Euclidean metric $\dS$. One could learn a more expressive metric
  (e.g., a Riemannian metric on $\SS$) from training data that assigns higher
  cost to contacts that are semantically significant (object boundaries) and
  lower cost to contacts that are semantically insignificant (texture boundaries).

\item \textbf{Efficient virtual substate enumeration}: A complete algorithmic
  specification of the virtual substate decomposition in SEBD, with complexity
  bounds for the virtual substate enumeration, is a necessary step toward a
  practical implementation.

\item \textbf{3D extension}: The current framework operates on 2D images. A
  natural extension to 3D volumetric images (CT, MRI, electron microscopy) would
  generalize the contact graph to 3D adjacency graphs and the holographic z-axis
  to a holographic ``temporal'' axis (the order of 3D contact resolutions).
\end{itemize}

\subsection{Closing Remarks}

The principle of individuation by negation is both ancient and novel. Its ancient
roots lie in the philosophical insight that identity is constituted by difference:
to be \emph{this} is to be \emph{not that}. Its novelty in the present work
lies in the precise mathematical articulation of this principle within the
framework of bounded resolvable spaces and its translation into a concrete,
efficient, theoretically principled algorithm for image analysis.

The S-Entropy Bidirectional Dijkstra is not merely a computational trick; it is
the algorithmic expression of a topological truth. Contact is invariant, irreducible,
and propagating. The holographic z-axis is not a geometric fiction; it is the
image's topological history, made spatial.

We believe that individuation by negation, as formalized here, offers a new
foundation for image analysis that is more philosophically coherent and
algorithmically principled than the dominant paradigm of edge detection by
gradient thresholding. The theoretical apparatus developed in this paper --- bounded
resolvable spaces, partition depth, contact topology, S-entropy coordinates,
and virtual substates --- provides the tools needed to explore this new foundation
systematically.

\printbibliography

\end{document}
